\chapter{Sauvegarde et restauration}
\label{ch:backup}

Une bibliothèque auto-hébergée n'a pas de filet de sécurité dans le
cloud. Ce chapitre couvre quoi sauvegarder, à quelle fréquence, et
comment restaurer.

\section{Quoi sauvegarder}

Trois artefacts forment ensemble une sauvegarde mybibli complète :

\begin{enumerate}
  \item \textbf{La base MariaDB.}\index{sauvegarde!base} Contient
        chaque titre, volume, lieu, contributeur, prêt,
        utilisateur, journal d'audit et paramètre. La perdre,
        c'est repartir de zéro.
  \item \textbf{Les images de couverture.}\index{sauvegarde!couvertures}
        Téléchargées au catalogage et persistées dans le volume
        Docker \texttt{mybibli\_covers} (ou dans le chemin hôte du
        bind mount si vous avez basculé selon
        section~\ref{sec:install-covers-bind}). Les perdre, le
        catalogue fonctionne toujours mais chaque page de titre
        affiche un placeholder générique tant que vous n'avez pas
        re-fetché les métadonnées titre par titre (issue~\#213 — les
        installs antérieures à 1.1.4 sans ce volume perdaient les
        couvertures à chaque mise à jour du conteneur).
  \item \textbf{Le fichier \texttt{.env}.}\index{sauvegarde!.env}
        Contient les mots de passe (base, clés d'API non encore
        migrées). Le perdre, un re-déploiement vous oblige à
        re-saisir chaque credential.
\end{enumerate}

Les images Docker elles-mêmes ne sont \emph{pas} critiques — vous
pouvez toujours les re-tirer depuis Docker Hub.

\section{Procédure de sauvegarde}

\subsection*{Base de données — MariaDB embarquée}

Depuis l'hôte qui exécute la pile compose :

\begin{lstlisting}
docker compose exec -T db mariadb-dump \
    -u root -p"$MYSQL_ROOT_PASSWORD" \
    --single-transaction --routines --triggers \
    mybibli \
    > backups/mybibli-$(date +%Y%m%d-%H%M%S).sql
\end{lstlisting}

Le drapeau \texttt{--single-transaction} est important : il fait un
snapshot cohérent sans verrouiller les tables pendant la durée du
dump. \texttt{--routines --triggers} préserve les routines stockées
(mybibli n'en utilise pas pour l'instant, mais le drapeau est peu
coûteux et future-proof).

Compressez le dump si l'espace disque est une préoccupation :

\begin{lstlisting}
gzip backups/mybibli-*.sql
\end{lstlisting}

\subsection*{Base de données — MariaDB externe}

Lancez \texttt{mariadb-dump} directement contre le serveur externe.
Les utilisateurs DSM Synology peuvent aussi planifier une tâche
\emph{Hyper Backup} qui cible le paquet MariaDB — elle produit une
sauvegarde binaire qui se restaure plus vite qu'un dump SQL.

\subsection*{Couvertures et \texttt{.env}}

Pour un répertoire de couvertures monté en bind-mount sur l'hôte,
une simple archive \texttt{tar} suffit :

\begin{lstlisting}
tar czf backups/mybibli-files-$(date +%Y%m%d).tar.gz \
    .env "$COVERS_HOST_PATH"
\end{lstlisting}

Pour l'install par défaut (volume Docker nommé), sauvegardez le
volume via Docker — le chemin hôte d'un volume nommé varie selon la
version de Docker et n'est pas stable, donc n'essayez pas de le
\texttt{tar} directement :

\begin{lstlisting}
docker run --rm \
    -v mybibli_covers:/source:ro \
    -v "$(pwd)/backups:/backup" \
    alpine \
    tar czf "/backup/mybibli-covers-$(date +%Y%m%d).tar.gz" -C /source .
tar czf "backups/mybibli-env-$(date +%Y%m%d).tar.gz" .env
\end{lstlisting}

Le conteneur \texttt{alpine} jetable monte le volume nommé en
lecture seule, le répertoire \texttt{backups/} hôte en écriture, et
écrit l'archive tar. Même forme pour restaurer — voir la section
Restauration plus bas.

\section{Cadence recommandée}

\begin{itemize}
  \item \textbf{Quotidien} — dump de base automatisé. Même sur un
        catalogue rarement modifié, un dump journalier garantit que
        la pire perte de données est bornée à une journée.
  \item \textbf{Hebdomadaire} — covers et \texttt{.env}. Ils ne
        changent que quand vous cataloguez de nouveaux titres ou
        faites tourner des clés ; hebdomadaire est amplement
        suffisant.
  \item \textbf{Copie hors-site} — au moins mensuelle. Rsync du
        répertoire \texttt{backups/} vers un NAS distant, un
        disque USB que vous échangez ou un bucket de stockage
        cloud. Un incendie qui consume le serveur domestique ne
        doit pas aussi consumer vos sauvegardes.
\end{itemize}

\begin{tip}
Une sauvegarde qui n'a jamais été restaurée n'est pas une
sauvegarde. Une fois par trimestre, restaurez votre dernier dump
dans un environnement Docker jetable et vérifiez que les compteurs
de titres et de prêts correspondent à ce que rapporte
\texttt{/admin > Health} en production.
\end{tip}

\section{Procédure de restauration}

\subsection*{Base de données}

Pour le setup avec base embarquée :

\begin{lstlisting}
docker compose stop mybibli
docker compose exec -T db mariadb -u root -p"$MYSQL_ROOT_PASSWORD" \
    -e "DROP DATABASE mybibli; CREATE DATABASE mybibli CHARACTER SET utf8mb4;"
zcat backups/mybibli-20260513-1200.sql.gz | \
    docker compose exec -T db mariadb -u root -p"$MYSQL_ROOT_PASSWORD" mybibli
docker compose start mybibli
\end{lstlisting}

Pour MariaDB externe, les mêmes commandes client \texttt{mariadb}
fonctionnent contre le serveur externe.

\subsection*{Couvertures et \texttt{.env}}

Pour un install bind-mount hôte :

\begin{lstlisting}
tar xzf backups/mybibli-files-20260513.tar.gz
\end{lstlisting}

Pour l'install par défaut (volume nommé), restaurez via un
conteneur jetable symétrique à la sauvegarde :

\begin{lstlisting}
docker compose stop mybibli
docker run --rm \
    -v mybibli_covers:/target \
    -v "$(pwd)/backups:/backup:ro" \
    alpine \
    sh -c "rm -rf /target/* && tar xzf /backup/mybibli-covers-20260513.tar.gz -C /target"
tar xzf backups/mybibli-env-20260513.tar.gz
docker compose start mybibli
\end{lstlisting}

Redémarrez la pile pour que le conteneur mybibli prenne en compte
le \texttt{.env} restauré.

\section{Migration vers un nouvel hôte}

La procédure de sauvegarde est aussi la procédure de migration :

\begin{enumerate}
  \item Sur l'ancien hôte : produire un dump de base et une archive
        tar frais.
  \item Transférer les deux sur le nouvel hôte.
  \item Sur le nouvel hôte : installer Docker, récupérer le fichier
        compose depuis la release correspondant à la version du
        dump, placer \texttt{.env} et \texttt{covers/} à côté,
        restaurer la base, démarrer la pile.
  \item Vérifier en comptant les titres sur \texttt{/admin >
        Health} et en contrôlant quelques pages de titre au
        hasard.
\end{enumerate}

\begin{warning}
Le dump est lié à la version de \emph{schéma}. Si l'hôte source
tournait en 1.1.0 et que le nouvel hôte tire 1.5.0, les migrations
du nouvel hôte mettront à niveau la base restaurée — c'est OK.
L'inverse (restaurer un dump 1.5.0 sur une installation 1.1.0)
n'est \emph{pas} supporté et échouera au boot. Restaurez toujours
sur un hôte exécutant la même version mybibli ou plus récente.
\end{warning}
